\documentclass[10pt,a4paper]{article}
\usepackage[margin=20mm,headheight=15pt,headsep=7mm,footskip=10mm]{geometry}
\usepackage{fontspec,amsmath,amssymb,mathtools,microtype}
\setmainfont{lmroman10-regular.otf}[BoldFont=lmroman10-bold.otf,ItalicFont=lmroman10-italic.otf,BoldItalicFont=lmroman10-bolditalic.otf]
\setsansfont{lmsans10-regular.otf}[BoldFont=lmsans10-bold.otf,ItalicFont=lmsans10-oblique.otf,BoldItalicFont=lmsans10-boldoblique.otf]
\setmonofont[Scale=.86]{lmmono10-regular.otf}[ItalicFont=lmmono10-italic.otf]
\usepackage{ragged2e}\usepackage{xcolor,booktabs,longtable,array,graphicx,enumitem,fancyhdr,needspace,titlesec,hyperref,xurl}
\definecolor{AtlasNavy}{HTML}{18334C}\definecolor{AtlasTeal}{HTML}{116D73}\definecolor{AtlasGray}{HTML}{52616D}
\hypersetup{colorlinks=true,linkcolor=AtlasTeal,urlcolor=AtlasTeal,pdfcreator={Complexity Atlas},unicode=true}
\urlstyle{same}\setlength{\parindent}{0pt}\setlength{\parskip}{4.8pt}\setlength{\emergencystretch}{3em}
\renewcommand{\baselinestretch}{1.08}\setcounter{secnumdepth}{0}\setlist{nosep,leftmargin=1.5em}
\titleformat{\section}{\Large\bfseries\sffamily\color{AtlasNavy}\raggedright}{}{0em}{}
\titleformat{\subsection}{\large\bfseries\sffamily\color{AtlasTeal}\raggedright}{}{0em}{}
\titleformat{\subsubsection}{\normalsize\bfseries\color{AtlasNavy}\raggedright}{}{0em}{}
\titlespacing*{\section}{0pt}{14pt}{6pt}\titlespacing*{\subsection}{0pt}{10pt}{5pt}
\pagestyle{fancy}\fancyhf{}\fancyhead[L]{\small\sffamily\color{AtlasGray}COMPLEXITY ATLAS}
\fancyhead[R]{\small\sffamily\color{AtlasGray}ENGLISH EDITION}\fancyfoot[C]{\small\thepage}\renewcommand{\headrulewidth}{.3pt}
\newcommand{\AtlasMath}[1]{\par\begingroup\centering\setbox0=\hbox{$\displaystyle #1$}\ifdim\wd0>1.4\linewidth\typeout{ATLAS WIDE FORMULA: \the\wd0}\fi\ifdim\wd0>\linewidth\resizebox{\linewidth}{!}{\box0}\else\box0\fi\par\endgroup}
\newcommand{\AtlasInlineMath}[1]{\begingroup\setbox0=\hbox{$#1$}\ifdim\wd0>.4\linewidth\endgroup\AtlasMath{#1}\else\endgroup$#1$\fi}
\pdfstringdefDisableCommands{\def\AtlasInlineMath#1{#1}}
\DeclareMathOperator{\polylog}{polylog}\DeclareMathOperator{\dist}{dist}\DeclareMathOperator{\diam}{diam}
\DeclareMathOperator{\rank}{rank}\DeclareMathOperator{\tr}{tr}\DeclareMathOperator{\diag}{diag}
\DeclareMathOperator{\prox}{prox}\DeclareMathOperator{\KL}{KL}\DeclareMathOperator{\Var}{Var}
\providecommand{\eps}{\varepsilon}\providecommand{\R}{\mathbb R}\providecommand{\E}{\mathbb E}
\providecommand{\mathbbm}[1]{\mathbf{#1}}\providecommand{\mathds}[1]{\mathbf{#1}}
\providecommand{\llbracket}{[\![}\providecommand{\rrbracket}{]\!]}
\begin{document}
\raggedright

{\LARGE\bfseries\sffamily\color{AtlasNavy} Finite-sum Lipschitz optimization with a strongly convex mean on a ball: exact selectable component values\par}\vspace{5pt}
\textbf{Setting ID.} \texttt{SO\_\allowbreak{}ZO\_\allowbreak{}E09c\_\allowbreak{}sc}

Single-objective optimization

Minimize \AtlasInlineMath{F} globally over \AtlasInlineMath{B_2(0,R)} and return a point with small objective gap.

\textbf{Status.} Other quantified gap present; Bounds missing

\textbf{TASK tags.} formulation: Global minimization; solution concept: Global minimization; task family: Single objective

\textbf{MODEL tags.} convexity: Average strong convexity; convexity: Strongly convex on feasible ball; smoothness: Component function lipschitz; structure: Finite sum

\textbf{ORACLE tags.} access order: Zeroth-order interface; interface family: Selected component value; public information: Public inputs explicitly declared; query domain: Full space

\textbf{ALGORITHMS tags.} adaptivity: Adaptive; randomization: Deterministic; randomization: Private randomness allowed; restriction: No additional span restriction explicitly imposed

\textbf{OUTPUT tags.} certificate: No additional output certificate required

\textbf{METRIC tags.} criterion: Objective gap

\textbf{GUARANTEE tags.} aggregation: Per instance; mode: Deterministic; mode: Probability

\textbf{RESOURCE tags.} budget: Fixed worst-case budget

\needspace{7\baselineskip}\section{1. Setting specification}
\needspace{4\baselineskip}\subsection{Task}
Minimize \AtlasInlineMath{F} globally over \AtlasInlineMath{B_2(0,R)} and return a point with small objective gap.

\needspace{4\baselineskip}\subsection{Admissible model}
\AtlasInlineMath{F(x)=n^{-1}\sum_{i=1}^n f_i(x)}, with \AtlasInlineMath{n\in\mathbb{N}}. Each component \AtlasInlineMath{f_i:\mathbb{R}^d\to\mathbb{R}} is globally convex.

Every component satisfies \AtlasInlineMath{|f_i(x)-f_i(y)|\le L\|x-y\|_2}; smoothness and a shared minimizer are not required. The optimization domain is \AtlasInlineMath{B_2(0,R)}, and boundary minimizers are allowed.

The mean objective \AtlasInlineMath{F} is required to be \AtlasInlineMath{\mu}-strongly convex only on \AtlasInlineMath{B_2(0,R)}, with budgets satisfying \AtlasInlineMath{0<\mu R\le L}. No additional strong-convexity assumption is imposed on the components or on the mean objective outside the ball.

\needspace{4\baselineskip}\subsection{Oracle interface}
The algorithm may actively choose \AtlasInlineMath{i\in\{1,\ldots,n\}} and \AtlasInlineMath{x\in\mathbb{R}^d}. The interface \AtlasInlineMath{(i,x)\mapsto f_i(x)} returns one exact component value. Components may be revisited; full-average values must be computed through component calls.

The public inputs are \AtlasInlineMath{d}, \AtlasInlineMath{n}, \AtlasInlineMath{L}, \AtlasInlineMath{R}, \AtlasInlineMath{\varepsilon}, \AtlasInlineMath{\mu}, together with the sets, kernels and auxiliary operations explicitly declared above. Information about the unknown objective and any undisclosed sample distribution is obtained only through the stated interface.

\needspace{4\baselineskip}\subsection{Output}
Return \AtlasInlineMath{\widehat{x}\in B_2(0,R)}. The output may be an unqueried point, an average or convex combination that remains in the stated space. No function value, gradient, inner solution, or accuracy certificate must accompany the output.

\needspace{4\baselineskip}\subsection{Loss or accuracy criterion}
The base loss is \AtlasInlineMath{\ell(I,x)=F(x)-F_*}, where \AtlasInlineMath{F_*=\inf_{u\in B_2(0,R)}F(u)}. This nonnegative loss is evaluated at the output, with tolerance \AtlasInlineMath{\varepsilon}.

\needspace{4\baselineskip}\subsection{Resource accounting}
One \AtlasInlineMath{(i,x)\mapsto f_i(x)} costs \AtlasInlineMath{1} component-value calls; recomputing the full average costs \AtlasInlineMath{n} calls. Every initialization, snapshot and differencing point is charged individually.

Only the declared oracle resource is charged. Computation and storage using existing information, together with explicitly known geometric and proximal operations, are free.

\needspace{4\baselineskip}\subsection{Algorithm views and success guarantees}
\needspace{4\baselineskip}\subsubsection{Deterministic}
Only deterministic strategies are allowed. Queries, stopping and postprocessing may be arbitrary measurable adaptive functions of the history. No linear-span or zero-respecting restriction is imposed.

Interactions are sequential, and every actual response in a batch is charged separately. Except for explicitly listed auxiliary interfaces, arithmetic, linear algebra, storage, and internal computation on available information are free, with no additional communication or memory restriction. Finite-difference steps must be finite and positive; real-number precision is not charged separately.

After fixing the public inputs, choose one algorithm that works for every admissible instance and reply law. The requirement \AtlasInlineMath{\ell(I,\widehat{x})\le \varepsilon} holds deterministically for each instance.

Oracle cost uses a deterministic worst-case budget: \AtlasInlineMath{N\le T} on every possible execution path. For multiple resources, coordinatewise caps apply jointly along the same path.

\needspace{4\baselineskip}\subsubsection{Randomized}
Private randomization is allowed, with the seed independent of the instance and external randomness. Queries, stopping and postprocessing may be arbitrary measurable adaptive functions of the history. No linear-span or zero-respecting restriction is imposed.

Interactions are sequential, and every actual response in a batch is charged separately. Except for explicitly listed auxiliary interfaces, arithmetic, linear algebra, storage, and internal computation on available information are free, with no additional communication or memory restriction. Finite-difference steps must be finite and positive; real-number precision is not charged separately.

After fixing the public inputs, choose one algorithm that works for every admissible instance and reply law. The error requirement is \AtlasInlineMath{\Pr(\ell(I,\widehat{x})\le \varepsilon)\ge2/3}.

Probability is over the algorithm's random seed only.

Oracle cost uses a deterministic worst-case budget: \AtlasInlineMath{N\le T} on every possible execution path. For multiple resources, coordinatewise caps apply jointly along the same path.

\textbf{Output qualification.} Random selection among previously queried points is permitted, subject to the declared success guarantee.

\needspace{7\baselineskip}\section{2. Notation and parameters}
\AtlasInlineMath{d}: Dimension of the variable space \textbf{Domain:} \AtlasInlineMath{d\in\mathbb{N},\ d\ge1}.

\AtlasInlineMath{n}: Number of finite-sum components \textbf{Domain:} \AtlasInlineMath{n\in\mathbb{N},\ n\ge1}.

\AtlasInlineMath{F=n^{-1}\sum_{i=1}^n f_i}

\AtlasInlineMath{L}: Public upper bound on the objective Lipschitz constant \textbf{Domain:} \AtlasInlineMath{L>0}.

\AtlasInlineMath{|f_i(x)-f_i(y)|\le L\|x-y\|_2} holds for every component.

\AtlasInlineMath{R}: Feasible-ball radius \textbf{Domain:} \AtlasInlineMath{R>0}.

The output must belong to \AtlasInlineMath{B_2(0,R)}; the query radius is specified separately in the interface.

\AtlasInlineMath{\varepsilon}: Error target; norm and squared-norm criteria use different thresholds \textbf{Domain:} \AtlasInlineMath{\varepsilon>0}.

\AtlasInlineMath{A}: Lipschitz accuracy ratio \textbf{Domain:} \AtlasInlineMath{A>0}.

\AtlasInlineMath{A=LR/\varepsilon}

\AtlasInlineMath{d_0}: Universal dimension threshold required by the lower bound \textbf{Domain:} \AtlasInlineMath{d_0>0}.

Determined by the applicable theorem conditions; retain a symbolic threshold when no numerical value is supplied.

\AtlasInlineMath{\varepsilon_Z}: Normalized accuracy threshold for the randomized exact-value lower bound \textbf{Domain:} \AtlasInlineMath{\varepsilon_Z>0}.

Determined by the applicable theorem conditions; retain a symbolic threshold when no numerical value is supplied.

\AtlasInlineMath{\mu}: Public lower bound on the strong-convexity modulus \textbf{Domain:} \AtlasInlineMath{\mu>0}.

\AtlasInlineMath{x,y,u}: Function arguments and dummy variables in the formulas \textbf{Domain:} \AtlasInlineMath{x,y,u\in\mathbb{R}^d}; spherical domains are embedded in \AtlasInlineMath{\mathbb{R}^{d+1}}.

Admissible ranges are specified separately by the task, model and oracle in this profile.

\AtlasInlineMath{\widehat{x}}: Final algorithm output \textbf{Domain:} \AtlasInlineMath{B_2(0,R)}.

\AtlasInlineMath{I}: Complete problem instance \textbf{Domain:} \AtlasInlineMath{I\in\mathcal{I}}.

Includes the unknown objective and its admissible feedback distribution or environment law. Public parameters are fixed before the quantifiers.

\AtlasInlineMath{\mathcal{I}}: Admissible instance class determined by the eight specification fields in this profile \textbf{Domain:} Instances satisfying the public parameter bounds.

\AtlasInlineMath{\mathcal{A}}: Algorithm or strategy \textbf{Domain:} Measurable adaptive strategies allowed in this profile.

The algorithm is selected before the hidden instance.

\AtlasInlineMath{\ell}: Base loss or residual \textbf{Domain:} \AtlasInlineMath{\ell:\mathcal{I}\times\mathcal{X}_{\mathrm{out}}\to[0,+\infty]}.

The exact formula appears in this profile's performance-criterion section.

\AtlasInlineMath{\mathbb{E},\Pr}: Expectation and probability \textbf{Domain:} Evaluated over the algorithmic and feedback randomness specified by this profile's guarantee.

Fix the instance first; aggregate cost and error separately.

\AtlasInlineMath{\inf,\sup}: Infimum and supremum \textbf{Domain:} Evaluated only over the variable sets explicitly specified in the formula.

In the risk definition, the infimum over algorithms precedes the supremum over instances or centers.

\AtlasInlineMath{\|\cdot\|_2,\langle\cdot,\cdot\rangle}: Euclidean norm and inner product \textbf{Domain:} \AtlasInlineMath{\|x\|_2=\sqrt{\langle x,x\rangle}}.

\AtlasInlineMath{\langle x,y\rangle=x^\top y}

\AtlasInlineMath{B_p(c,R)}: Closed norm ball \textbf{Domain:} \AtlasInlineMath{p\in[1,\infty],\ R\ge0}.

\AtlasInlineMath{B_p(c,R)=\{x:\|x-c\|_p\le R\}}; \AtlasInlineMath{B_2} denotes the Euclidean ball.

\AtlasInlineMath{N,T}: Actual call count and declared budget \textbf{Domain:} \AtlasInlineMath{N,T\in\mathbb{N}_0}.

A fixed budget bounds every execution path; the Markov specification uses an expected budget.

\AtlasInlineMath{O,\Omega,\Theta,\widetilde{O}}: Asymptotic upper bounds, lower bounds, matching orders, and upper bounds with hidden logarithms \textbf{Domain:} Within the joint parameter window specified for each complexity regime.

Dependence of hidden constants on fixed parameters is stated separately for each bound. Tilde notation may hide logarithmic factors only in the declared variables.

\AtlasInlineMath{c,C,c_\alpha,C_\alpha}: Positive constant in a proof or asymptotic bound \textbf{Domain:} \AtlasInlineMath{0<c,C,c_\alpha,C_\alpha<\infty}.

Unsubscripted constants have the dependencies stated in the bound; constants subscripted by \AtlasInlineMath{\alpha} may depend only on the fixed smoothness order.

\AtlasInlineMath{\log,\log_+,e}: Natural logarithm, truncated logarithm, and Euler's number \textbf{Domain:} \AtlasInlineMath{z>0}.

\AtlasInlineMath{\log_+(z)=\max\{0,\log z\}}; \AtlasInlineMath{e=\exp(1)}, and \AtlasInlineMath{\log_2} denotes the base-two logarithm.

\AtlasInlineMath{f_i,i}: Finite-sum components and their indices \textbf{Domain:} \AtlasInlineMath{f_i:\mathbb{R}^d\to\mathbb{R},\ i\in\{1,\ldots,n\}}.

Components are unknown; indices may be selected actively and revisited at a cost.

\AtlasInlineMath{F}: Finite-sum mean objective \textbf{Domain:} \AtlasInlineMath{F:\mathbb{R}^d\to\mathbb{R}}.

\AtlasInlineMath{F=n^{-1}\sum_{i=1}^n f_i}

\AtlasInlineMath{F_*}: Optimal value of the finite-sum objective over the optimization domain \textbf{Domain:} \AtlasInlineMath{F_*\in\mathbb{R}}.

\AtlasInlineMath{F_*=\inf_{x\in\mathcal{X}}F(x)}.

\AtlasInlineMath{\mathcal{X}_{\mathrm{fun}},\mathcal{X}_{\mathrm{out}},\mathcal{X}}: Function domain, admissible output domain and optimization domain \textbf{Domain:} Each is the Euclidean space or subset explicitly specified in this profile.

\AtlasInlineMath{\mathcal{X}_{\mathrm{out}}=B_2(0,R)} and \AtlasInlineMath{\mathcal{X}=B_2(0,R)}; the function domain is the region on which the model imposes regularity.

\AtlasInlineMath{N^*}: Optimal worst-case oracle complexity satisfying the complete guarantee \textbf{Domain:} \AtlasInlineMath{N^*\in[0,+\infty]}.

Take the infimum over admissible algorithms, followed by the worst case over the instances of this profile. With multiple resources, use this scalar only for an explicitly specified resource projection.

\needspace{12\baselineskip}\section{3. Complexity bracket}
\needspace{4\baselineskip}\subsection{Compared quantities}
\needspace{4\baselineskip}\subsubsection{Q1 - target accuracy cost}
\textbf{View.} Deterministic

\AtlasMath{C^\star_{v,\pi}(\varepsilon;\theta)=\inf\{B\ge0:\exists\mathcal A\in\mathfrak A_v,\ C_\pi(\mathcal A,I)\le B\ \text{and}\ G_v(\mathcal A,I;\varepsilon)\ \forall I\}}
\textbf{Definition.} one exact selected component scalar value

\textbf{Scope.} declared\_joint\_parameter\_window

\textbf{Optimization direction.} minimize

Here v is the selected algorithm view, G\_v its complete success guarantee, I an admissible fixed instance, theta the remaining parameters, and C\_pi the declared resource projection. C\_pi uses the view's fixed-path or expected-work budget convention, including initialization. No change of guarantee or scalarization is implicit.

\needspace{4\baselineskip}\subsubsection{Q2 - target accuracy cost}
\textbf{View.} Randomized

\AtlasMath{C^\star_{v,\pi}(\varepsilon;\theta)=\inf\{B\ge0:\exists\mathcal A\in\mathfrak A_v,\ C_\pi(\mathcal A,I)\le B\ \text{and}\ G_v(\mathcal A,I;\varepsilon)\ \forall I\}}
\textbf{Definition.} one exact selected component scalar value

\textbf{Scope.} declared\_joint\_parameter\_window

\textbf{Optimization direction.} minimize

Here v is the selected algorithm view, G\_v its complete success guarantee, I an admissible fixed instance, theta the remaining parameters, and C\_pi the declared resource projection. C\_pi uses the view's fixed-path or expected-work budget convention, including initialization. No change of guarantee or scalarization is implicit.

\needspace{4\baselineskip}\subsection{R1 - Strongly convex finite sums on a ball: comparable window}
\textbf{Region.} \AtlasInlineMath{0<\mu R\le L/2}, \AtlasInlineMath{0<\varepsilon\le\min\{LR/8,LRd/[512n(d+\log(16n))]\}}.

\textbf{Views.} Deterministic

\textbf{Quantity.} Q1

\textbf{Resource.} one exact selected component scalar value

\textbf{Status.} Other quantified gap

\textbf{Lower bound.} B2

\AtlasMath{\Omega(nd)}
\textbf{Upper bound.} B1

\AtlasMath{O\!\left(n\min\{dA^2\log(edA),dL^2\log(edA)/(\mu\varepsilon),d^2\log(ed)\log(4A)\}\right)}
\textbf{Gap or comparison.} \AtlasInlineMath{O(\min\{A^2\log(edA),L^2\log(edA)/(\mu\varepsilon),d\log(ed)\log(4A)\})}

\needspace{4\baselineskip}\subsection{R2 - Strongly convex finite sums on a ball: a compatible lower bound remains to be verified}
\textbf{Region.} \AtlasInlineMath{0<\mu R\le L} and \AtlasInlineMath{0<\varepsilon\le LR/8}. In addition, require \AtlasInlineMath{(\mu R>L/2\ \text{or}\ \varepsilon>LRd/[512n(d+\log(16n))])} and \AtlasInlineMath{\varepsilon<\min\{LR,(L-\mu R)^2/(2\mu)\}}.

\textbf{Views.} Deterministic

\textbf{Quantity.} Q1

\textbf{Resource.} one exact selected component scalar value

\textbf{Status.} Lower bound missing

\textbf{Lower bound.} No admitted bound for this comparison

\textbf{Upper bound.} B1

\AtlasMath{O\!\left(n\min\{dA^2\log(edA),dL^2\log(edA)/(\mu\varepsilon),d^2\log(ed)\log(4A)\}\right)}
\needspace{4\baselineskip}\subsection{R3 - Strongly convex finite sums on a ball: zero-query regime}
\textbf{Region.} \AtlasInlineMath{\varepsilon\ge\min\{LR,(L-\mu R)^2/(2\mu)\}}. When \AtlasInlineMath{\mu R=L}, every positive target accuracy lies in this regime.

\textbf{Views.} Deterministic; Randomized

\textbf{Quantity.} Q1; Q2

\textbf{Resource.} one exact selected component scalar value

\textbf{Status.} Optimal

\textbf{Lower bound.} B4, B9

\AtlasMath{0}
\textbf{Upper bound.} B3, B8

\AtlasMath{0}
\needspace{4\baselineskip}\subsection{R4 - Strongly convex finite sums on a ball: comparable window}
\textbf{Region.} \AtlasInlineMath{0<\mu R\le L/2}, \AtlasInlineMath{0<\varepsilon\le\min\{LR/8,LRd/[512n(d+\log(16n))]\}}.

\textbf{Views.} Randomized

\textbf{Quantity.} Q2

\textbf{Resource.} one exact selected component scalar value

\textbf{Status.} Other quantified gap

\textbf{Lower bound.} B7

\AtlasMath{\Omega(nd)}
\textbf{Upper bound.} B5, B6

\AtlasMath{O(\min\{dA^2,dL^2/(\mu\varepsilon),nd^2\log(ed)\log(4A)\})}
\textbf{Gap or comparison.} \AtlasInlineMath{O(\min\{A^2/n,L^2/(n\mu\varepsilon),d\log(ed)\log(4A)\})}

\needspace{4\baselineskip}\subsection{R5 - Strongly convex finite sums on a ball: a compatible lower bound remains to be verified}
\textbf{Region.} \AtlasInlineMath{0<\mu R\le L} and \AtlasInlineMath{0<\varepsilon\le LR/8}. In addition, require \AtlasInlineMath{(\mu R>L/2\ \text{or}\ \varepsilon>LRd/[512n(d+\log(16n))])} and \AtlasInlineMath{\varepsilon<\min\{LR,(L-\mu R)^2/(2\mu)\}}.

\textbf{Views.} Randomized

\textbf{Quantity.} Q2

\textbf{Resource.} one exact selected component scalar value

\textbf{Status.} Lower bound missing

\textbf{Lower bound.} No admitted bound for this comparison

\textbf{Upper bound.} B5, B6

\AtlasMath{O(\min\{dA^2,dL^2/(\mu\varepsilon),nd^2\log(ed)\log(4A)\})}
\needspace{7\baselineskip}\section{4. Bounds and supporting evidence}
The following conditions apply to every bound in this profile, in addition to each bound's stated region and algorithm scope.

\textbf{Oracle mapping.} Each source component value corresponds to one selectable component query; computing the exact mean requires accessing every component.

\textbf{Guarantee mapping.} Error uses this profile's accuracy scale and aggregation order. The source theorem's expectation, probability, and asymptotic quantifiers are retained separately.

\textbf{Output mapping.} The output need only lie in this setting's prescribed space. Permitted averages and unqueried outputs must still satisfy its feasibility conditions.

\textbf{Resource mapping.} Atomic calls under a fixed cap on every execution path. All feedback obtained during initialization, differencing and inner iterations is charged.

\needspace{4\baselineskip}\subsection{B1 - Two-point descent with component sampling, or deterministic localization}
\textbf{Admission.} Included in the bracket

\textbf{Direction.} upper

\textbf{View.} Deterministic

\textbf{Witness or lower-bound scope.} Deterministic witness

\AtlasMath{O\!\left(n\min\{dA^2\log(edA),dL^2\log(edA)/(\mu\varepsilon),d^2\log(ed)\log(4A)\}\right)}
\textbf{Sources.} \href{https://arxiv.org/html/2607.16558v2}{[1]} \href{https://www.mathnet.ru/php/archive.phtml?jrnid=mzm&option_lang=eng&paperid=1697&wshow=paper}{[2]} \href{https://arxiv.org/html/2610.00345v1}{[3]} \href{https://arxiv.org/abs/2109.08858}{[4]}

\textbf{Conditions.} \AtlasInlineMath{0<\mu R\le L,\ 0<\varepsilon\le LR/8}.

\textbf{Local argument.} A1, A4

\needspace{4\baselineskip}\subsection{B2 - Adaptive-measurement lower bound for independent Gaussian components}
\textbf{Admission.} Included in the bracket

\textbf{Direction.} lower

\textbf{View.} Deterministic

\textbf{Witness or lower-bound scope.} All randomized adaptive algorithms

\AtlasMath{\Omega(nd)}
\textbf{Conditions.} \AtlasInlineMath{0<\mu\le L/(2R),\ 0<\varepsilon\le LRd/[512n(d+\log(16n))]}.

\textbf{Local argument.} A3

\needspace{4\baselineskip}\subsection{B3 - Return the public initial point}
\textbf{Admission.} Included in the bracket

\textbf{Direction.} upper

\textbf{View.} Deterministic

\textbf{Witness or lower-bound scope.} Deterministic witness

\AtlasMath{0}
\textbf{Conditions.} \AtlasInlineMath{\varepsilon\ge\min\{LR,(L-\mu R)^2/(2\mu)\}}.

\textbf{Local argument.} A5, A6

\needspace{4\baselineskip}\subsection{B4 - Nonnegative oracle cost}
\textbf{Admission.} Included in the bracket

\textbf{Direction.} lower

\textbf{View.} Deterministic

\textbf{Witness or lower-bound scope.} All randomized adaptive algorithms

\AtlasMath{0}
\textbf{Conditions.} \AtlasInlineMath{\varepsilon\ge\min\{LR,(L-\mu R)^2/(2\mu)\}}.

\textbf{Local argument.} A5, A6

\needspace{4\baselineskip}\subsection{B5 - Two-point descent with component sampling, or deterministic localization}
\textbf{Admission.} Included in the bracket

\textbf{Direction.} upper

\textbf{View.} Randomized

\textbf{Witness or lower-bound scope.} Randomized witness

\AtlasMath{O(\min\{dA^2,dL^2/(\mu\varepsilon),nd^2\log(ed)\log(4A)\})}
\textbf{Sources.} \href{https://jmlr.org/papers/volume18/16-632/16-632.pdf}{[5]} \href{https://arxiv.org/html/2607.16558v2}{[1]} \href{https://www.mathnet.ru/php/archive.phtml?jrnid=mzm&option_lang=eng&paperid=1697&wshow=paper}{[2]} \href{https://arxiv.org/html/2610.00345v1}{[3]} \href{https://arxiv.org/abs/2109.08858}{[4]}

\textbf{Conditions.} \AtlasInlineMath{0<\mu R\le L,\ 0<\varepsilon\le LR/8}; success probability at least 2/3.

\textbf{Local argument.} A8, A4

\needspace{4\baselineskip}\subsection{B6 - Two-point descent with component sampling, or deterministic localization}
\textbf{Admission.} Included in the bracket

\textbf{Direction.} upper

\textbf{View.} Randomized

\textbf{Witness or lower-bound scope.} Deterministic witness

\AtlasMath{O\!\left(n\min\{dA^2\log(edA),dL^2\log(edA)/(\mu\varepsilon),d^2\log(ed)\log(4A)\}\right)}
\textbf{Sources.} \href{https://arxiv.org/html/2607.16558v2}{[1]} \href{https://www.mathnet.ru/php/archive.phtml?jrnid=mzm&option_lang=eng&paperid=1697&wshow=paper}{[2]} \href{https://arxiv.org/html/2610.00345v1}{[3]} \href{https://arxiv.org/abs/2109.08858}{[4]}

\textbf{Conditions.} \AtlasInlineMath{0<\mu R\le L,\ 0<\varepsilon\le LR/8}.

\textbf{Local argument.} A1, A4

\needspace{4\baselineskip}\subsection{B7 - Adaptive-measurement lower bound for independent Gaussian components}
\textbf{Admission.} Included in the bracket

\textbf{Direction.} lower

\textbf{View.} Randomized

\textbf{Witness or lower-bound scope.} All randomized adaptive algorithms

\AtlasMath{\Omega(nd)}
\textbf{Conditions.} \AtlasInlineMath{0<\mu\le L/(2R),\ 0<\varepsilon\le LRd/[512n(d+\log(16n))]}.

\textbf{Local argument.} A3

\needspace{4\baselineskip}\subsection{B8 - Return the public initial point}
\textbf{Admission.} Included in the bracket

\textbf{Direction.} upper

\textbf{View.} Randomized

\textbf{Witness or lower-bound scope.} Deterministic witness

\AtlasMath{0}
\textbf{Conditions.} \AtlasInlineMath{\varepsilon\ge\min\{LR,(L-\mu R)^2/(2\mu)\}}.

\textbf{Local argument.} A5, A6

\needspace{4\baselineskip}\subsection{B9 - Nonnegative oracle cost}
\textbf{Admission.} Included in the bracket

\textbf{Direction.} lower

\textbf{View.} Randomized

\textbf{Witness or lower-bound scope.} All randomized adaptive algorithms

\AtlasMath{0}
\textbf{Conditions.} \AtlasInlineMath{\varepsilon\ge\min\{LR,(L-\mu R)^2/(2\mu)\}}.

\textbf{Local argument.} A5, A6

\needspace{7\baselineskip}\section{5. Local arguments}
\needspace{4\baselineskip}\subsection{A1 - Deterministic extraction of an approximate subgradient}
\textbf{Scope.} Exact real scalar values; finite positive probe scales; free real arithmetic and no bit-precision charge. Ball-restricted variants keep all probes inside the ball; global variants allow all-space probes.

A deterministic procedure obtains an \AtlasInlineMath{\eta}-accurate subgradient at a point within distance \AtlasInlineMath{\zeta} of the requested point using at most \AtlasInlineMath{7d\lceil\log_2(2L\sqrt d/\eta)\rceil} values. It yields the convex upper \AtlasInlineMath{O(dA^2\log(edA))} and the strongly convex upper \AtlasInlineMath{O(dL^2\log(edA)/(\mu\varepsilon))}, with \AtlasInlineMath{A=LR/\varepsilon}.

For a convex one-dimensional restriction, maintain an active interval whose every subderivative lies in \AtlasInlineMath{[\ell,u]}. Divide it into six equal pieces and query its seven grid points. Their secant slopes satisfy \AtlasInlineMath{s_0\le\cdots\le s_5}. The intervals \AtlasInlineMath{[q_1,q_2]} and \AtlasInlineMath{[q_4,q_5]} have protected subgradient envelopes \AtlasInlineMath{[s_0,s_2]} and \AtlasInlineMath{[s_3,s_5]}. Since \AtlasInlineMath{(s_2-s_0)+(s_5-s_3)\le u-\ell}, selecting the narrower one halves the envelope width and divides the interval length by six.

Set \AtlasInlineMath{\tau=\eta/\sqrt d}, \AtlasInlineMath{k=\lceil\log_2(2L/\tau)\rceil}, and \AtlasInlineMath{r_1=\zeta/2}. At coordinate \AtlasInlineMath{j}, run \AtlasInlineMath{k} subdivisions from half-width \AtlasInlineMath{r_j}, move to the final midpoint, and retain the two protected secants of denominator \AtlasInlineMath{h_j=2r_j6^{-k}}. Set \AtlasInlineMath{r_{j+1}=\tau h_j/(16L)}. Then \AtlasInlineMath{r_{j+1}/r_j\le1/8}, and all later displacement is at most \AtlasInlineMath{\tau h_j/(8L)}. Translating a protected secant by this displacement changes its slope by at most \AtlasInlineMath{\tau/4}. Thus every subgradient of the final common point \AtlasInlineMath{y} has coordinate \AtlasInlineMath{j} in \AtlasInlineMath{[\ell_j-\tau/4,u_j+\tau/4]}. The midpoint vector differs from every such subgradient by at most \AtlasInlineMath{3\eta/4} in Euclidean norm. All probe and protection points lie within \AtlasInlineMath{\zeta} of the original point. The scales may be very small, but are positive and require the same finite number of queries.

For a globally Lipschitz convex function and a public projected set \AtlasInlineMath{K\subseteq B_2(R)} containing zero, the extracted vector satisfies

\AtlasMath{f(z)\ge f(x)+\langle\widehat g,z-x\rangle-2L\zeta-2R\eta\quad(z\in K).}

Choose \AtlasInlineMath{\zeta=\varepsilon/(16L)}, \AtlasInlineMath{\eta=\varepsilon/(16R)}, so the model error is \AtlasInlineMath{\varepsilon/4} and \AtlasInlineMath{\|\widehat g\|\le2L} in the stated nontrivial range. Projected subgradient descent with step \AtlasInlineMath{R/(2L\sqrt N)} gives averaged gap at most \AtlasInlineMath{2LR/\sqrt N+\varepsilon/4}. Taking \AtlasInlineMath{N\ge64A^2} proves the convex rate. The probes may leave \AtlasInlineMath{K} only in the all-space variant.

For a function defined on the ball, keep iterates in \AtlasInlineMath{B_2(R-2\delta)}. For strong convexity, apply extraction to \AtlasInlineMath{h(x)=f(x)-\mu\|x\|^2/2}, which is convex and at most \AtlasInlineMath{2L}-Lipschitz because \AtlasInlineMath{\mu R\le L}. Choose \AtlasInlineMath{\delta=\gamma/(16L)}, vector error \AtlasInlineMath{\eta=\gamma/(8R)}, and add \AtlasInlineMath{\mu x} back. Then \AtlasInlineMath{\|\widehat g\|\le4L} and

\AtlasMath{f(z)\ge f(x)+\langle\widehat g,z-x\rangle+\mu\|z-x\|^2/2-\gamma.}

With \AtlasInlineMath{a_t=2/[\mu(t+1)]}, projection, and \AtlasInlineMath{D_t=\|x_t-x_*\|^2},

\AtlasMath{f(x_t)-f(x_*)-\gamma\le\frac{\mu(t-1)}4D_t-\frac{\mu(t+1)}4D_{t+1}+\frac{16L^2}{\mu(t+1)}.}

Multiply by \AtlasInlineMath{t} and sum. The initial distance coefficient is zero, and the remaining distance terms telescope. The weighted average has gap at most \AtlasInlineMath{32L^2/[\mu(T+1)]+\gamma}. Taking \AtlasInlineMath{\gamma=\varepsilon/8} and \AtlasInlineMath{T+1\ge128L^2/(\mu\varepsilon)}, including shrinkage loss at most \AtlasInlineMath{2L\delta}, proves the result. Ordinary convex projected descent gives the \AtlasInlineMath{A^2} branch. Deterministic value localization supplies the separate \AtlasInlineMath{d^2\log(ed)\log(4A)} branch; take the minimum only among applicable branches.

For a known body \AtlasInlineMath{B_2(r_{\rm in})\subseteq K\subseteq B_2(R_{\rm out})}, the identity \AtlasInlineMath{(1-a)K+aB_2(r_{\rm in})\subseteq K} supplies feasible probe room. Choose \AtlasInlineMath{a=\gamma/(16LR_{\rm out})} and probe radius at most \AtlasInlineMath{ar_{\rm in}/2}. Shrinkage costs at most \AtlasInlineMath{LaR_{\rm out}}. The query count is independent of the probe radius and uses \AtlasInlineMath{R_{\rm out}} in place of \AtlasInlineMath{R}. A finite-sum mean value costs \AtlasInlineMath{n} component values, multiplying these deterministic bounds by \AtlasInlineMath{n}.

\textbf{Limits of the argument.} Local argument and specification transfer; cited external theorems remain primary dependencies.

\textbf{Sources.} \href{https://arxiv.org/html/2607.16558v2}{[1]} \href{https://www.mathnet.ru/php/archive.phtml?jrnid=mzm&option_lang=eng&paperid=1697&wshow=paper}{[2]}

\needspace{4\baselineskip}\subsection{A2 - Exact recovery at equal smoothness and strong convexity}
\textbf{Scope.} Global \AtlasInlineMath{\mu}-strong convexity and \AtlasInlineMath{\beta}-smoothness with \AtlasInlineMath{\mu=\beta}; exact scalar values and unrestricted finite queries.

The objective is \AtlasInlineMath{f(x)=\beta\|x\|^2/2-b^\top x+c}. Values at \AtlasInlineMath{0,e_1,\ldots,e_d} recover a minimizer in \AtlasInlineMath{d+1} calls.

Strong monotonicity, Cauchy--Schwarz, and Lipschitz continuity of the gradient imply

\AtlasMath{\beta\|x-y\|^2\le\langle\nabla f(x)-\nabla f(y),x-y\rangle\le\|\nabla f(x)-\nabla f(y)\|\,\|x-y\|\le\beta\|x-y\|^2.}

Equality throughout forces \AtlasInlineMath{\nabla f(x)-\nabla f(y)=\beta(x-y)}. Therefore \AtlasInlineMath{b_j=\beta/2+f(0)-f(e_j)} and \AtlasInlineMath{x_*=b/\beta}. For a finite-sum average, each of these \AtlasInlineMath{d+1} average values costs \AtlasInlineMath{n} component values, giving \AtlasInlineMath{n(d+1)} calls. A one-component specialization has \AtlasInlineMath{n=1}. Restricted query domains require an admissible nonzero coordinate scale in place of unit vectors.

\textbf{Limits of the argument.} Local argument and specification transfer; cited external theorems remain primary dependencies.

\needspace{4\baselineskip}\subsection{A3 - An adaptive \AtlasInlineMath{nd} lower bound for selected-component values}
\textbf{Scope.} Fixed scalar-call cap; arbitrary measurable adaptive randomized queries, including unbounded query points; only the explicitly stated average or component localization promises.

For smooth finite sums, \AtlasInlineMath{T>nd/2} is necessary at \AtlasInlineMath{\varepsilon\le\mu R^2/512}. For globally Lipschitz components the same conclusion holds at

\AtlasMath{\varepsilon\le\frac{LRd}{512n[d+\log(16n)]}.}

The displayed constants give success probability \AtlasInlineMath{2/3}; an expected-error transfer uses the same statement at \AtlasInlineMath{3\varepsilon}.

Take independent \AtlasInlineMath{a_i\sim N(0,s^2I_d)}. A component query, after subtracting a known common term, reveals one linear measurement of \AtlasInlineMath{a_i}. Conditional on the random seed and the transcript, the components remain independent Gaussian vectors with covariance \AtlasInlineMath{s^2P_i}, where \AtlasInlineMath{P_i} projects onto the complement of the span queried in that component. Hence the mean \AtlasInlineMath{\bar a} has covariance \AtlasInlineMath{\Sigma=(s^2/n^2)\sum_iP_i} and

\AtlasMath{\operatorname{tr}\Sigma\ge s^2(nd-T)/n^2,\qquad\Sigma\preceq(s^2/n)I.}

For any Gaussian vector \AtlasInlineMath{Z}, any proposed center \AtlasInlineMath{b}, and \AtlasInlineMath{t>0}, Gaussian integration gives \AtlasInlineMath{\mathbb E e^{-t\|Z-b\|^2}\le(1+2t\operatorname{tr}\Sigma)^{-1/2}}. Taking \AtlasInlineMath{t=1/(2r^2)} yields

\AtlasMath{\Pr(\|Z-b\|^2\le r^2)\le\sqrt{\frac{er^2}{r^2+\operatorname{tr}\Sigma}}.}

In particular, \AtlasInlineMath{r^2\le\operatorname{tr}\Sigma/8} makes this at most \AtlasInlineMath{\sqrt e/3}.

For the smooth family set \AtlasInlineMath{f_i(x)=\mu\|x\|^2/2-\mu\langle a_i,x\rangle} and \AtlasInlineMath{s^2=nR^2/(16d)}. The mean is minimized at \AtlasInlineMath{\bar a}. The event \AtlasInlineMath{E=\{\|\bar a\|\le R\}} has probability at least \AtlasInlineMath{15/16}. For \AtlasInlineMath{T\le nd/2}, \AtlasInlineMath{\operatorname{tr}\Sigma\ge R^2/32}. Objective accuracy \AtlasInlineMath{\varepsilon} means squared center error at most \AtlasInlineMath{2\varepsilon/\mu}, which is at most \AtlasInlineMath{\operatorname{tr}\Sigma/8} on the displayed window. The unconditional success probability is at most \AtlasInlineMath{\sqrt e/3}; conditioning on \AtlasInlineMath{E} only after this calculation bounds valid-instance success by \AtlasInlineMath{16\sqrt e/45<2/3}. Choose \AtlasInlineMath{\mu=\beta} for the merely convex smooth class, or the declared \AtlasInlineMath{\mu\le\beta} for the strongly convex class. Individual component centers are not bounded in this construction.

For globally Lipschitz components define the radial Huber function

\AtlasMath{h_R(x)=\begin{cases}\|x\|^2/2,&\|x\|\le R,\\R\|x\|-R^2/2,&\|x\|>R,\end{cases}\quad q=L/(2R),\quad f_i=q(h_R-\langle a_i,x\rangle).}

Choose \AtlasInlineMath{s^2=R^2/[8(d+\log(16n))]}. A Gaussian tail bound and union bound give \AtlasInlineMath{\Pr(\max_i\|a_i\|\le R)\ge15/16}. On this event each component is globally convex and \AtlasInlineMath{L}-Lipschitz. Inside the feasible ball the mean gap is \AtlasInlineMath{q\|x-\bar a\|^2/2}. The same posterior argument gives \AtlasInlineMath{\operatorname{tr}\Sigma\ge R^2d/[16n(d+\log(16n))]}, proving the stated window. Every unbounded query still reveals just one linear measurement because the common Huber term is known. Inside the ball this family has curvature \AtlasInlineMath{qI}, so it also proves the average-strongly-convex finite-sum lower for every public \AtlasInlineMath{\mu\le L/(2R)}.

These are joint \AtlasInlineMath{nd} arguments. Separate \AtlasInlineMath{\Omega(n)} and \AtlasInlineMath{\Omega(d)} statements would justify only a maximum, not their product.

\textbf{Limits of the argument.} Local argument and specification transfer; cited external theorems remain primary dependencies.

\needspace{4\baselineskip}\subsection{A4 - Finite-sum simulation and single-component inheritance}
\textbf{Scope.} The average objective, component regularity, public geometry, algorithm permissions, and guarantee are retained exactly.

One average value costs \AtlasInlineMath{n} selected-component values. A lower bound obtained by making all components equal transfers with cost one. Bounds copied from a one-component model without an \AtlasInlineMath{n} factor apply only on \AtlasInlineMath{n=1}.

For \AtlasInlineMath{F=n^{-1}\sum_i f_i}, evaluate every component at the same point and average. This exactly simulates any whole-objective value, including finite-difference probes outside the feasible ball when all-space queries are permitted. Thus deterministic convex acceleration and localization give \AtlasInlineMath{O(nd\sqrt Q)} and \AtlasInlineMath{O(nd^2\log(e+Q))}; strongly convex acceleration, restart, and localization give

\AtlasMath{O\!\left(nd\min\{\sqrt Q,\sqrt\kappa[1+\log_+(Q/\kappa)],d\log(e+Q)\}\right).}

The near-unit-condition-number method uses the midpoint-curvature contraction \AtlasInlineMath{(\kappa-1)/(\kappa+1)} with a finite-difference error small enough for the prescribed horizon, and each value again costs \AtlasInlineMath{n}. At \AtlasInlineMath{\kappa=1}, use exact recovery of the average quadratic. For a nonsmooth mean, multiply deterministic localization or approximate-subgradient simulation by \AtlasInlineMath{n}.

For lower bounds set every component equal to a legal single-objective hard instance. A requested component index adds no information and one hard-oracle value supplies its response. This transfers the source lower bound without multiplying by \AtlasInlineMath{n}, with exactly its own dimension, accuracy, and condition-number windows. Combine it with the independent joint \AtlasInlineMath{nd} lower only where both apply. If a row imports a one-component upper or endpoint formula without an \AtlasInlineMath{n} multiplier, its condition is \AtlasInlineMath{n=1}.

ARCS Corollaries 6 and 8 supply the separate randomized finite-sum upper branches. The public ball's linear minimization oracle is free in this model; the global minimizer lies in that ball, so the initial potential is bounded using \AtlasInlineMath{\beta R^2} (for the convex branch, \AtlasInlineMath{D_0\le7\beta R^2}). Its finite-difference component queries are legal in the all-space domain. A constant-factor tighter expected nonnegative gap gives fixed success by Markov's inequality. These algorithms remain randomized and do not enter a deterministic view.

\textbf{Limits of the argument.} Local argument and specification transfer; cited external theorems remain primary dependencies.

\textbf{Sources.} \href{https://arxiv.org/html/2610.00345v1}{[3]} \href{https://arxiv.org/abs/2109.08858}{[4]} \href{https://arxiv.org/html/2607.16558v2}{[1]}

\needspace{4\baselineskip}\subsection{A5 - Zero-query certificates}
\textbf{Scope.} The stated sufficient accuracy regions only; no assertion that every certificate is sharp.

A public feasible point whose objective gap is at most \AtlasInlineMath{G_0} solves every target \AtlasInlineMath{\varepsilon\ge G_0} with exactly zero calls.

For a convex mean objective that is \AtlasInlineMath{L}-Lipschitz on a ball of radius \AtlasInlineMath{R}, the center has gap at most \AtlasInlineMath{LR}. A sample Lipschitz modulus \AtlasInlineMath{G} with \AtlasInlineMath{\mathbb E G^r\le L^r}, \AtlasInlineMath{r>1}, implies \AtlasInlineMath{\mathbb E G\le L}, so the same certificate applies to moment models. Integrability at one point ensures that the mean is defined.

If a globally \AtlasInlineMath{\beta}-smooth objective has a global minimizer within distance \AtlasInlineMath{R} of the public initial point, its gradient vanishes at that minimizer and the smoothness inequality gives gap at most \AtlasInlineMath{\beta R^2/2}. In a declared norm, use the corresponding radius in that norm. For \AtlasInlineMath{B_p(R)} with Euclidean regularity this radius is \AtlasInlineMath{R} when \AtlasInlineMath{p\le2} and \AtlasInlineMath{R d^{1/2-1/p}} when \AtlasInlineMath{p\ge2}.

For a known positive semidefinite quadratic matrix \AtlasInlineMath{A}, the center has gap at most \AtlasInlineMath{\lambda_{\max}(A)R^2/2}; if \AtlasInlineMath{A=0}, the loss is identically zero. A public initial-gap promise \AtlasInlineMath{\Delta} directly gives the certificate \AtlasInlineMath{G_0=\Delta}. For the Markov potential \AtlasInlineMath{r_0}, its companion initial iterate has gap at most \AtlasInlineMath{\mu r_0}. These are different certificates with different model assumptions; in particular, an unknown matrix with no public spectral bound does not acquire a uniform certificate.

Call counts are nonnegative. Together with the stated algorithm, this proves exact zero complexity on each certified region. The upper/lower ratio is undefined, so it is stored as null rather than as \AtlasInlineMath{0/0}.

\textbf{Limits of the argument.} Local argument and specification transfer; cited external theorems remain primary dependencies.

\needspace{4\baselineskip}\subsection{A6 - The center certificate for a strongly convex Lipschitz function}
\textbf{Scope.} A convex function on \AtlasInlineMath{B_2(R)}, \AtlasInlineMath{\mu}-strongly convex and \AtlasInlineMath{L}-Lipschitz; interior subgradients are taken without an artificial boundary normal cone.

The class is nonempty only if \AtlasInlineMath{\mu R\le L}, and

\AtlasMath{f(0)-\min_{B_2(R)} f\le\min\{LR,(L-\mu R)^2/(2\mu)\}.}

In particular, the center is optimal when \AtlasInlineMath{\mu R=L}.

Let \AtlasInlineMath{g_0\in\partial f(0)}. If \AtlasInlineMath{g_0\ne0}, put \AtlasInlineMath{u=g_0/\|g_0\|} and choose \AtlasInlineMath{g_t\in\partial f(tu)} for \AtlasInlineMath{0<t<R}. Strong monotonicity gives \AtlasInlineMath{\langle g_t,u\rangle\ge\|g_0\|+\mu t}. Interior subgradients have norm at most \AtlasInlineMath{L}, hence \AtlasInlineMath{\|g_0\|\le L-\mu R} by letting \AtlasInlineMath{t\uparrow R}. If \AtlasInlineMath{g_0=0}, the same inequality along any direction gives the feasibility condition \AtlasInlineMath{\mu R\le L}.

Strong convexity implies \AtlasInlineMath{f(z)\ge f(0)+\langle g_0,z\rangle+\mu\|z\|^2/2\ge f(0)-\|g_0\|^2/(2\mu)}. Combine this with the independent Lipschitz center bound \AtlasInlineMath{LR}. The argument uses only the mean function, so it also applies when a finite-sum mean is strongly convex only on the feasible ball.

\textbf{Limits of the argument.} Local argument and specification transfer; cited external theorems remain primary dependencies.

\needspace{4\baselineskip}\subsection{A7 - Ball smoothing and the same-sample directional moment}
\textbf{Scope.} A fresh sample each round; the same sample supplies both values; sample Lipschitz modulus \AtlasInlineMath{G} on the stated query neighborhood; \AtlasInlineMath{\mathbb E G^r\le L^r}, \AtlasInlineMath{1<r\le2}.

With ball smoothing \AtlasInlineMath{f_h(x)=\mathbb E_{U\sim\operatorname{Unif}(B_2)}f(x+hU)} and a sphere direction \AtlasInlineMath{u},

\AtlasMath{g_h(x)=\frac d{2h}[F(x+hu;Z)-F(x-hu;Z)]u,\quad\mathbb E g_h=\nabla f_h,\quad\mathbb E\|g_h\|^r\le(\sqrt{2d}L)^r.}

For \AtlasInlineMath{r=2}, the bound is \AtlasInlineMath{2dL^2}.

For each fixed sample, the scalar difference \AtlasInlineMath{a(u)=F(x+hu;Z)-F(x-hu;Z)} is odd on the sphere, hence has mean zero, and has Lipschitz constant at most \AtlasInlineMath{2hG}. Sphere Poincare gives \AtlasInlineMath{\mathbb E_u a(u)^2\le4h^2G^2/(d-1)} for \AtlasInlineMath{d\ge2}. Thus \AtlasInlineMath{\mathbb E_u\|g_h\|^2\le d^2G^2/(d-1)\le2dG^2}; for \AtlasInlineMath{d=1} the difference quotient is directly bounded by \AtlasInlineMath{G}. Apply Jensen only to the direction expectation and then integrate over the sample:

\AtlasMath{\mathbb E_{Z,u}\|g_h\|^r\le\mathbb E_Z(\mathbb E_u\|g_h\|^2)^{r/2}\le(2d)^{r/2}\mathbb E G^r.}

No second moment of \AtlasInlineMath{G} is needed when \AtlasInlineMath{r<2}. The divergence theorem for ball convolution gives the unbiasedness identity. Integrability follows from the Lipschitz difference bound. Convexity gives \AtlasInlineMath{f_h\ge f}, and \AtlasInlineMath{|f_h-f|\le hL}; mean strong convexity is preserved on the region where every translated point satisfies it. All statements hold conditionally on the past when the next sample and direction are fresh. Reusing an old adaptively explored sample does not automatically preserve this conditional moment calculation.

\textbf{Limits of the argument.} Local argument and specification transfer; cited external theorems remain primary dependencies.

\textbf{Sources.} \href{https://jmlr.org/papers/volume18/16-632/16-632.pdf}{[5]}

\needspace{4\baselineskip}\subsection{A8 - Projected optimization with same-sample two-point values}
\textbf{Scope.} Optimization/output in \AtlasInlineMath{B_2(R)}; all two-point queries remain in the prescribed ball; mean strong convexity suffices. A uniformly sampled finite-sum component is also a valid fresh sample.

For a \AtlasInlineMath{\mu}-strongly-convex mean with \AtlasInlineMath{\mathbb E G^2\le L^2}, \AtlasInlineMath{K} rounds and \AtlasInlineMath{T=2K} scalar calls give

\AtlasMath{\mathbb E[f(\bar x_K)-f_*]\le\frac{4dL^2}{\mu(K+1)}+2hL.}

The merely convex branch uses \AtlasInlineMath{O(d(LR/\varepsilon)^2)} values.

Run the paired estimator on \AtlasInlineMath{K_h=B_2(R-h)} and let \AtlasInlineMath{x_{h,*}} minimize \AtlasInlineMath{f_h} there. The shrinkage comparator \AtlasInlineMath{y=(1-h/R)x_*} gives \AtlasInlineMath{f_h(y)\le f_*+2hL}. Set \AtlasInlineMath{V=2dL^2}, \AtlasInlineMath{\eta_t=2/[\mu(t+1)]}, and \AtlasInlineMath{D_t=\mathbb E\|x_t-x_{h,*}\|^2}. Projection and strong convexity yield

\AtlasMath{\mathbb E[f_h(x_t)-f_h(x_{h,*})]\le\frac{\mu(t-1)}4D_t-\frac{\mu(t+1)}4D_{t+1}+\frac{V}{\mu(t+1)}.}

Multiply by \AtlasInlineMath{t}, telescope, and average with weights \AtlasInlineMath{2t/[K(K+1)]}. The initial distance coefficient is zero and the last is nonpositive, giving \AtlasInlineMath{2V/[\mu(K+1)]}. Add the comparator error. For expected accuracy choose, for example, \AtlasInlineMath{h=\varepsilon/(8L)} and \AtlasInlineMath{K+1\ge8dL^2/(\mu\varepsilon)}. For success probability \AtlasInlineMath{2/3}, use \AtlasInlineMath{h=\varepsilon/(12L)} and \AtlasInlineMath{K+1\ge24dL^2/(\mu\varepsilon)}, making the expected nonnegative gap at most \AtlasInlineMath{\varepsilon/3} before Markov.

For convex objectives, the projected averaging recursion with step \AtlasInlineMath{R/\sqrt{VK}} gives expected optimization error at most \AtlasInlineMath{R\sqrt{V/K}}, plus the same smoothing/shrinkage error. Uniform component sampling makes the estimator unbiased for the finite-sum mean and costs exactly two selected-component values per round. Neither the oracle nor the components need be stochastic in that case; the algorithm supplies the randomness. These bounds also hold under an almost-sure Lipschitz promise, since it implies the required second moment. The Gaussian lower construction in the separate lower lemma does not acquire that stronger promise.

\textbf{Limits of the argument.} Local argument and specification transfer; cited external theorems remain primary dependencies.

\needspace{7\baselineskip}\section{References}
[1] Haihan Zhang, Chenheng Zhang, Zhiquan Qi and Zhouchen Lin. Near-Optimal Lower Bounds for Randomized Algorithms in Exact Value Zeroth-Order Convex Optimization. arXiv:2607.16558v2, 2026-07-30.

\url{https://arxiv.org/html/2607.16558v2}

[2] V. Yu. Protasov. Algorithms for approximate calculation of the minimum of a convex function from its values. Mathematical Notes 59(1):69--74, 1996.

\url{https://www.mathnet.ru/php/archive.phtml?jrnid=mzm&option_lang=eng&paperid=1697&wshow=paper}

[3] Dorn, Dvinskikh, Loginov, Osmanov, Shestakov, Tupitsa; Gasnikov. Lower Bounds For Gradient-Free Convex Optimization And Convex-Concave Saddle-Point Problems. arXiv:2610.00345v1.

\url{https://arxiv.org/html/2610.00345v1}

[4] Xiyuan Wei
Bin Gu
Heng Huang. An Accelerated Variance-Reduced Conditional Gradient Sliding Algorithm for First-order and Zeroth-order Optimization. 2021. arXiv2109.08858v1 (2021-09-18).

\url{https://arxiv.org/abs/2109.08858}

[5] Ohad Shamir. An Optimal Algorithm for Bandit and Zero-Order Convex Optimization with Two-Point Feedback. JMLR 18(52):1--11, 2017.

\url{https://jmlr.org/papers/volume18/16-632/16-632.pdf}

\end{document}
